\section{Statistical implications}

The comparisons in \cref{obj:thm:point-frontier,obj:thm:interval-frontier} distinguish sampling uncertainty from the additional cost of recovering outcome means across baseline cells. In the benchmark of \cref{obj:def:rate}, sampling contributes \(n^{-1}\) to squared error, while the dimension contribution is weighted by the square of the arrival deficit \(1-q\). Accordingly, the minimax root mean squared error has order
\[
n^{-1/2}+(1-q)\min\left\{1,\frac{d}{n\ell_n}\right\},
\qquad \ell_n=\log(e+n).
\]
The comparison constants are uniform over positive integers \(n,d\) and known arrival floors \(q\in[1/2,1]\). This uniformity permits sample size, baseline dimension, and outcome arrival to vary jointly.

Before saturation, when \(d\le n\ell_n\), the dimension contribution to root mean squared error is \((1-q)d/(n\ell_n)\). It matches the sampling scale when
\[
(1-q)d=\sqrt n\,\ell_n.
\]
Below this balance, sampling determines the order of precision; above it, the arrival-weighted dimension term determines that order. For a fixed arrival floor below one, the transition therefore occurs at baseline dimension of order \(\sqrt n\,\ell_n/(1-q)\), provided this lies within the unsaturated range. As the arrival floor increases, a larger baseline dimension becomes compatible with sampling-scale precision. In particular, \(1-q\le n^{-1/2}\) guarantees precision of order \(n^{-1/2}\) uniformly over all positive baseline-support bounds. Complete arrival gives the exact endpoint described in \cref{obj:prop:complete-arrival-reduction}.

Saturation occurs at \(d\ge n\ell_n\), where the dimension contribution to squared error equals \((1-q)^2\). Increasing the support bound within this regime leaves the benchmark unchanged. With a fixed arrival floor below one, the minimax risk then remains bounded away from zero as sample size grows along such a sequence. An arrival floor approaching one instead reduces the saturated contribution directly. Thus the rate separates the effect of increasing sample size relative to dimension from the effect of improving outcome arrival.

More generally, let \(d_n\) be a sequence of positive baseline-support bounds and \(q_n\in[1/2,1]\) a sequence of known arrival floors, with \(n\to\infty\). The point-estimation comparison implies that minimax squared risk tends to zero precisely when
\[
(1-q_n)\min\left\{1,\frac{d_n}{n\ell_n}\right\}\longrightarrow0.
\]
For an arrival floor fixed below one, this condition becomes \(d_n=o(n\ell_n)\). For \(q_n\to1\), it is satisfied for every dimension sequence. Sampling-order squared risk, of order \(n^{-1}\), is attained precisely when the displayed dimension scale is \(O(n^{-1/2})\). These statements describe the joint sequences through the same finite-sample comparison, with constants independent of the sequence.

Honest inference translates these regimes into uniform coverage with corresponding expected precision. For each fixed \(\alpha\in(0,1/2)\), \cref{obj:thm:interval-frontier} gives minimax worst-law expected length of order
\[
n^{-1/2}+(1-q)\min\left\{1,\frac{d}{n\ell_n}\right\},
\]
with constants depending only on \(\alpha\). The interval criterion in \cref{obj:def:interval-class,obj:def:length-risk} requires coverage at least \(1-\alpha\) for every law in the specified class at each sample size. Consequently, the joint-sequence condition above also characterizes when this optimal worst-law expected length tends to zero. Sampling-dominated regimes yield honest expected length of order \(n^{-1/2}\), while saturation yields length of order \(n^{-1/2}+1-q\). The lower comparison concerns the worst-law expected length among honest connected intervals; the construction in \cref{obj:def:mm-interval} attains the corresponding order throughout the class.

The converse also has a useful implication for model enlargement. The sampling and dimension lower bounds in \cref{obj:lem:parametric-floor-infrastructure,obj:thm:normalized-converse,obj:lem:prior-reduction} are supported by randomized submodels with a constant surrogate. Any broader law class on the same observed experiment, with the same treatment-effect target, that contains these submodels inherits their lower bounds: its worst-law criterion includes the laws used in the converse. This establishes a necessary precision scale for such containing classes. The matching upper bounds characterize attainable precision over the finite-cell randomized-MAR model in \cref{obj:def:law-class}, with its binary measurements, known arrival floor, and baseline-support bound. Together, the two directions identify how sampling, dimension, and outcome arrival govern optimal precision in that model.
